\documentclass[11pt]{article}

\usepackage[a4paper,margin=1in]{geometry}
\usepackage{amsmath}
\usepackage{booktabs}
\usepackage{float}
\usepackage{graphicx}
\usepackage{hyperref}
\usepackage{listings}
\usepackage{xcolor}
\usepackage{enumitem}
\usepackage{caption}
\emergencystretch=3em
\sloppy

\definecolor{codebg}{RGB}{245,245,245}
\lstset{
  backgroundcolor=\color{codebg},
  basicstyle=\ttfamily\small,
  breaklines=true,
  frame=single,
  language=Python,
  columns=fullflexible,
  keepspaces=true,
  showstringspaces=false,
  numbers=left,
  numberstyle=\tiny,
  xleftmargin=1.5em
}

\title{Simulation and Performance Analysis of a Direct-Mapped\\
Instruction and Data Cache}
\author{Sudarshan Sudhakar \\ Roll No: CS23B2007}
\date{\today}

\begin{document}
\maketitle

\begin{abstract}
This report describes the design and evaluation of a Python-based cache
simulator that models a direct-mapped Instruction Cache (I-Cache) and
Data Cache (D-Cache) for a system with a 64K-word main memory, a 2K-word
cache, a 16-word block size, and 32-bit words/addresses. Three
representative memory access patterns -- sequential, repeated, and
strided -- are simulated for both instruction and data streams. For each
run the simulator reports total accesses, cache lookups, hits, misses,
hit rate, and miss rate, and further decomposes every miss into
Compulsory (Cold), Conflict, and Capacity categories using the
Hill--Smith three-C classification method. The results are used to
discuss how spatial and temporal locality, together with the
direct-mapped placement policy, govern cache efficiency.
\end{abstract}

\tableofcontents

% =====================================================================
\section{System Specification}
% =====================================================================

\begin{table}[H]
\centering
\begin{tabular}{ll}
\toprule
Parameter & Value \\
\midrule
Main memory size       & 64K words ($2^{16}$ words) \\
Cache size (each of I-Cache, D-Cache) & 2K words ($2^{11}$ words) \\
Block size              & 16 words ($2^4$ words) \\
Word size               & 32 bits \\
Address size            & 32 bits \\
Mapping scheme          & Direct Mapped \\
\bottomrule
\end{tabular}
\caption{Given system parameters}
\end{table}

\subsection{Derived Cache Geometry}

Number of lines (blocks) in the cache:
\[
\text{Lines} = \frac{\text{Cache size}}{\text{Block size}} = \frac{2048}{16} = 128 \text{ lines}
\]

Since the cache is direct mapped, a 32-bit word address is split into
three fields -- \emph{Tag}, \emph{Index}, and \emph{Offset}:

\[
\text{Offset bits} = \log_2(16) = 4, \qquad
\text{Index bits} = \log_2(128) = 7, \qquad
\text{Tag bits} = 32 - 7 - 4 = 21
\]

\begin{table}[H]
\centering
\begin{tabular}{l p{3.2cm} p{3.2cm} p{3.6cm}}
\toprule
Field & Tag & Index & Offset \\
\midrule
Width (bits) & 21 & 7 & 4 \\
Purpose & identifies which memory block & selects the 1-of-128 cache line & selects the word within a 16-word block \\
\bottomrule
\end{tabular}
\caption{Address decomposition used by the simulator (\texttt{cache.decompose})}
\end{table}

\textbf{Assumption:} the assignment specifies a single ``2K-word cache''
but asks that both an I-Cache and a D-Cache be modelled. We interpret
this as a Harvard-style organisation in which the I-Cache and the
D-Cache are two \emph{independent} caches, each built to the given
specification (2K words, direct mapped, 16-word blocks). They are
simulated completely independently -- an instruction access can never
hit or evict a line in the D-Cache and vice versa.

% =====================================================================
\section{Simulator Design}
% =====================================================================

The simulator is written in Python 3 and organised into three modules:

\begin{itemize}[leftmargin=1.6em]
  \item \texttt{cache.py} -- the direct-mapped cache model, the address
        decomposition logic, and the three-C (Cold/Conflict/Capacity)
        miss classifier.
  \item \texttt{patterns.py} -- generators that produce the address
        streams for the three required access patterns, for
        instructions and for data.
  \item \texttt{simulator.py} -- the driver that runs every
        pattern$\times$\{instruction, data\} combination through a
        fresh cache instance, prints a formatted report, and dumps the
        results to \texttt{results/results.json}.
\end{itemize}

\subsection{Direct-Mapped Cache Model}

Each cache is represented by two parallel arrays of length 128 (one
entry per line): a \texttt{valid} bit and a \texttt{tag}. On every
access the address is decomposed into (tag, index, offset); the
\texttt{index} selects a single candidate line, and \emph{exactly one}
tag comparison (one ``lookup'') is performed:

\begin{itemize}[leftmargin=1.6em]
  \item If the line is valid and its stored tag matches $\Rightarrow$ \textbf{hit}.
  \item Otherwise $\Rightarrow$ \textbf{miss}; the block is loaded into
        that line (allocate-on-miss), overwriting whatever was
        previously resident there.
\end{itemize}

Because a direct-mapped cache checks only one line per access, the
number of \emph{cache lookups} reported below is always equal to the
number of accesses (one tag comparison per access) -- this is
structurally different from a set-associative cache, where a single
access can require several parallel tag comparisons.

\subsection{Classifying Misses into Cold / Conflict / Capacity}

We use the standard Hill--Smith ``Three-C'' method, implemented in
\texttt{ThreeCClassifier}:

\begin{enumerate}[leftmargin=1.6em]
  \item \textbf{Cold (Compulsory) miss} -- the referenced memory block
        has never been accessed before in this run. Such a miss would
        occur even in an infinite cache and is unavoidable on first
        touch.
  \item A reference \textbf{fully-associative LRU cache} of the
        \emph{same capacity} (128 blocks) is run in parallel, purely
        as a measuring stick.
        \begin{itemize}
          \item If the direct-mapped cache misses, the block is not
                cold, but the reference LRU cache \emph{also} misses on
                it $\Rightarrow$ \textbf{Capacity miss} (the working set
                genuinely exceeds the cache's total capacity, so no
                placement policy could have avoided this miss).
          \item If the direct-mapped cache misses but the reference LRU
                cache \emph{hits} $\Rightarrow$ \textbf{Conflict miss}
                (the miss is an artefact of direct-mapped placement --
                two blocks competing for the same index evicted each
                other, even though the overall working set would have
                fit).
        \end{itemize}
\end{enumerate}

Every miss produced by the direct-mapped cache is assigned to exactly
one of these three buckets, and the simulator asserts
$\text{Cold}+\text{Conflict}+\text{Capacity} = \text{Total misses}$
after every run as a self-check.

\subsection{Access Patterns Simulated}

All addresses are 32-bit \emph{word} addresses into the 64K-word memory.
Instruction streams are generated in the region starting at word~0,
data streams in a disjoint region starting at word $0{\times}8000$
(32768), purely so the two are easy to reason about independently; the
caches themselves are simulated completely separately regardless.

\begin{enumerate}[leftmargin=1.6em]
\item \textbf{Sequential (spatial locality).} 4096 consecutive word
      addresses (256 cache blocks, i.e. 2$\times$ the cache's 128-line
      capacity), modelling straight-line instruction execution or a
      linear array scan with no reuse.

\item \textbf{Repeated (temporal locality).} Two sub-cases:
      \begin{itemize}
        \item \emph{Fits in cache}: a 32-word instruction loop body
              (2 blocks) executed 200 times, and a 16-word data
              window (1 block) reused 500 times -- both comfortably
              smaller than the 128-line cache.
        \item \emph{Exceeds cache capacity}: a 3072-word (192-block,
              1.5$\times$ capacity) instruction working set and a
              4096-word (256-block, 2$\times$ capacity) data working
              set, each swept 4 times, to show what happens when a
              ``hot loop'' is actually too large to stay resident.
      \end{itemize}

\item \textbf{Strided.} Fixed-stride sweeps with stride
      $\in\{2,4,16\}$ words over 1024 accesses, showing spatial
      locality eroding as the stride grows (stride 16 = one full
      block, i.e. zero reuse within a block). A fourth, dedicated
      \emph{conflict-demonstration} case round-robins across 4
      addresses spaced exactly one cache-size (2048 words) apart, 256
      times -- these four blocks alias to the very same cache index,
      which isolates \textbf{conflict misses} from capacity effects
      (a fully-associative cache of equal size would need only 4 of
      its 128 lines to hold this entire working set).
\end{enumerate}

% =====================================================================
\section{Results}
% =====================================================================

All figures below are the direct, un-edited output of
\texttt{simulator.py} (see \texttt{results/results.json}).

\subsection{Sequential Access Pattern}
\begin{table}[H]
\centering
\resizebox{\textwidth}{!}{%
\begin{tabular}{lrrrrrrrrr}
\toprule
Cache / Case & Acc. & Lookups & Hits & Misses & Hit\% & Miss\% & Cold & Conflict & Capacity \\
\midrule
I-Cache (sequential) & 4096 & 4096 & 3840 & 256 & 93.75\% & 6.25\% & 256 & 0 & 0 \\
D-Cache (sequential) & 4096 & 4096 & 3840 & 256 & 93.75\% & 6.25\% & 256 & 0 & 0 \\
\bottomrule
\end{tabular}%
}
\caption{Sequential access pattern -- I-Cache and D-Cache results}
\label{tab:sequential}
\end{table}

\begin{table}[H]
\centering
\resizebox{\textwidth}{!}{%
\begin{tabular}{lrrrrrrrrr}
\toprule
Cache / Case & Acc. & Lookups & Hits & Misses & Hit\% & Miss\% & Cold & Conflict & Capacity \\
\midrule
I-Cache (working set fits in cache) & 6400 & 6400 & 6398 & 2 & 99.97\% & 0.03\% & 2 & 0 & 0 \\
D-Cache (working set fits in cache) & 8000 & 8000 & 7999 & 1 & 99.99\% & 0.01\% & 1 & 0 & 0 \\
I-Cache (working set $>$ cache capacity) & 12288 & 12288 & 11712 & 576 & 95.31\% & 4.69\% & 192 & 0 & 384 \\
D-Cache (working set $>$ cache capacity) & 16384 & 16384 & 15360 & 1024 & 93.75\% & 6.25\% & 256 & 0 & 768 \\
\bottomrule
\end{tabular}%
}
\caption{Repeated (temporal locality) access pattern -- I-Cache and D-Cache results}
\label{tab:repeated}
\end{table}

\begin{table}[H]
\centering
\resizebox{\textwidth}{!}{%
\begin{tabular}{lrrrrrrrrr}
\toprule
Cache / Case & Acc. & Lookups & Hits & Misses & Hit\% & Miss\% & Cold & Conflict & Capacity \\
\midrule
I-Cache (stride = 2) & 1024 & 1024 & 896 & 128 & 87.50\% & 12.50\% & 128 & 0 & 0 \\
D-Cache (stride = 2) & 1024 & 1024 & 896 & 128 & 87.50\% & 12.50\% & 128 & 0 & 0 \\
I-Cache (stride = 4) & 1024 & 1024 & 768 & 256 & 75.00\% & 25.00\% & 256 & 0 & 0 \\
D-Cache (stride = 4) & 1024 & 1024 & 768 & 256 & 75.00\% & 25.00\% & 256 & 0 & 0 \\
I-Cache (stride = 16) & 1024 & 1024 & 0 & 1024 & 0.00\% & 100.00\% & 1024 & 0 & 0 \\
D-Cache (stride = 16) & 1024 & 1024 & 0 & 1024 & 0.00\% & 100.00\% & 1024 & 0 & 0 \\
\bottomrule
\end{tabular}%
}
\caption{Strided access pattern (stride = 2, 4, 16 words) -- I-Cache and D-Cache results}
\label{tab:strided}
\end{table}

\begin{table}[H]
\centering
\resizebox{\textwidth}{!}{%
\begin{tabular}{lrrrrrrrrr}
\toprule
Cache / Case & Acc. & Lookups & Hits & Misses & Hit\% & Miss\% & Cold & Conflict & Capacity \\
\midrule
I-Cache (index-aliasing conflict demo) & 1024 & 1024 & 0 & 1024 & 0.00\% & 100.00\% & 4 & 1020 & 0 \\
D-Cache (index-aliasing conflict demo) & 1024 & 1024 & 0 & 1024 & 0.00\% & 100.00\% & 4 & 1020 & 0 \\
\bottomrule
\end{tabular}%
}
\caption{Strided access, conflict-demonstration case -- 4 addresses exactly one cache-size (2048 words) apart, round-robined 256 times}
\label{tab:conflict-demo}
\end{table}


% tables.tex emits 4 \begin{table}...\end{table} blocks in order:
% sequential, repeated, strided(2/4/16), strided conflict-demo.
% Section headers below annotate where each one lands only in prose.

\subsection*{Repeated Access Pattern}
See Table~\ref{tab:repeated} above.

\subsection*{Strided Access Pattern}
See Tables~\ref{tab:strided} and~\ref{tab:conflict-demo} above.

% =====================================================================
\section{Analysis and Discussion}
% =====================================================================

\subsection{Sequential Access (Table~\ref{tab:sequential})}

Both caches show an identical \textbf{93.75\% hit rate}. With a 16-word
block, every block delivers 1 miss followed by 15 hits as long as the
stream never revisits an evicted block -- exactly $256/4096 = 6.25\%$
miss rate. Every single miss is classified as \textbf{Cold}; there is
no conflict or capacity component because each of the 256 blocks
touched is accessed exactly once and then abandoned, so nothing is ever
evicted and re-requested. This is the textbook signature of
\emph{spatial locality}: large blocks amortise the fixed cost of a
miss over many subsequent hits, and the hit rate is entirely
determined by the block size ($\frac{15}{16}=93.75\%$), independent of
cache size or associativity.

\subsection{Repeated Access (Table~\ref{tab:repeated})}

\textbf{Working set within capacity.} With only 2 (I) or 1 (D) blocks
being swept repeatedly, hit rates reach \textbf{99.97\%} and
\textbf{99.99\%}: the entire working set is loaded once (a handful of
cold misses) and then serviced from the cache for hundreds of
iterations. This is the purest demonstration of \emph{temporal
locality} -- exactly the loop / loop-carried-variable scenario the
pattern is meant to model.

\textbf{Working set exceeding capacity.} Once the swept region grows
past the cache's 128-line capacity (192 blocks for I, 256 blocks for
D), the hit rate falls to 95.31\% / 93.75\% and, critically, the miss
breakdown now contains a large \textbf{Capacity} component (384 and
768 misses respectively) on top of the unavoidable cold misses -- with
\emph{zero} conflict misses. Because the sweep revisits blocks in the
same round-robin order every pass, a fully-associative LRU cache of
equal size thrashes in exactly the same way as the direct-mapped
cache, so every non-cold miss here is attributable purely to
insufficient capacity, not to index collisions. This isolates the
effect the assignment asks about: a loop or array that is individually
``hot'' but simply too large to fit will thrash regardless of how
clever the placement policy is.

\subsection{Strided Access (Tables~\ref{tab:strided}, \ref{tab:conflict-demo})}

\textbf{Increasing stride erodes spatial locality.} For a plain
strided sweep with no revisits, the hit rate falls monotonically as
stride increases: $87.50\%\to75.00\%\to0.00\%$ for stride
$2\to4\to16$. With a 16-word block, a stride of 2 still lands 8 of the
16 words in each block before moving to the next block (7 hits after
the first miss on the block, so with count 1024 that is
$\tfrac{1}{2}\cdot 1{,}024$ block-touches down to 128 unique blocks,
each contributing 1 cold miss and 7 hits); stride 4 similarly halves
the reuse again; and once the stride equals the block size (16), every
single access lands in a brand-new block, spatial locality is
completely destroyed, and the miss rate hits \textbf{100\%}. As with
the sequential pattern, every miss here is classified as Cold --
nothing is ever revisited, so there is nothing for a
conflict/capacity distinction to act on.

\textbf{Isolating conflict misses.} The dedicated
conflict-demonstration case (Table~\ref{tab:conflict-demo}) makes the
cost of direct mapping explicit: only 4 blocks are ever touched (4
cold misses), and this tiny working set would trivially fit in
\emph{any} cache with $\geq 4$ lines of \emph{any} associativity. Yet
because all 4 addresses are exactly one cache size (2048 words) apart,
they map to the \emph{same} index and continuously evict one another.
The result is a \textbf{100\% miss rate} with 1020 of the 1024 misses
classified as pure \textbf{Conflict} misses -- the single largest
source of avoidable misses observed in this study, and one that a
capacity argument alone cannot explain (the reference fully-associative
cache of equal total size hits almost every time). This is precisely
the scenario direct-mapped caches are known to handle poorly, and is
the reason set-associative caches exist.

\subsection{Instruction Cache vs.\ Data Cache}

Across every pattern, the I-Cache and D-Cache results are structurally
identical in shape (same hit/miss ratios, same 3C breakdown), because
both caches use the same geometry and are driven by address streams
constructed with the same access pattern (only the base address
differs). This is an expected artefact of testing the two caches with
synthetic streams; in a real workload the I-Cache would additionally
be shaped by branch behaviour (taken/not-taken changes sequential
fetch into effectively strided/discontinuous access) while the D-Cache
would be shaped by the specific data structures being manipulated
(arrays $\to$ sequential/strided, scalars and loop counters $\to$
repeated). The simulator's independent I/D modelling reflects a
Harvard-style split cache, where instruction fetch pressure can never
evict useful data, and vice versa -- avoiding the interference that
would occur in a shared unified cache.

\subsection{Locality of Reference and Cache Mapping -- Summary}

\begin{itemize}[leftmargin=1.6em]
\item \textbf{Spatial locality} (exploited by block size) is the
      dominant factor for sequential and low-stride access: hit rate
      is essentially $1 - \frac{1}{\text{words per block used}}$, and
      degrades sharply once the access stride approaches or exceeds
      the block size.
\item \textbf{Temporal locality} (exploited by cache capacity and
      residency) is the dominant factor for repeated/loop-style
      access, \emph{provided the working set fits inside the cache}.
      When it doesn't, capacity misses appear regardless of placement
      policy.
\item \textbf{Cache mapping policy} (direct-mapped vs.\ associative)
      only becomes visible once addresses \emph{alias} onto the same
      index. The conflict-demonstration experiment shows this can be
      catastrophic (100\% miss rate) even for a working set that is
      trivially small relative to total cache capacity -- the defining
      weakness of direct mapping relative to set-associative or
      fully-associative organisations.
\end{itemize}

% =====================================================================
\section{Conclusion}
% =====================================================================

The simulator confirms, with concrete measurements, the standard
qualitative predictions of cache theory for a direct-mapped
organisation: sequential access performance is governed by block size
(spatial locality), repeated access performance is governed by whether
the working set fits within cache capacity (temporal locality), and
strided access exposes both a gradual loss of spatial locality as
stride grows and, in the worst case (addresses aliasing to the same
index), severe conflict misses that no amount of spare cache capacity
can prevent. Because the I-Cache and D-Cache are modelled as
independent direct-mapped structures, instruction and data traffic
never interfere with one another, isolating each pattern's effect
cleanly. These results motivate the two classic remedies covered in a
computer organisation course: larger/well-chosen block sizes to
capture more spatial locality, and higher associativity (or software
techniques such as array padding/blocking) to eliminate pathological
conflict misses like the one demonstrated in
Table~\ref{tab:conflict-demo}.

% =====================================================================
\appendix
\section{Source Code}
% =====================================================================

The complete, runnable source (\texttt{cache.py}, \texttt{patterns.py},
\texttt{simulator.py}, \texttt{generate\_report\_tables.py}) is
submitted alongside this report in the \texttt{src/} directory, and the
raw JSON results in \texttt{results/results.json}. The core cache
model and 3C classifier are reproduced below for reference.

\lstinputlisting[caption={\texttt{cache.py} -- direct-mapped cache model and 3C classifier}]{../src/cache.py}

\end{document}
